\usepackage[utf8]{inputenc}
\usepackage[T1]{fontenc}
\usepackage[spanish]{babel}
\decimalpoint % separador decimal con punto, como en las tablas de la fuente original
\usepackage{lmodern,microtype}
\usepackage{amsmath,amssymb,amsfonts}
\usepackage{geometry}
\geometry{margin=2.54cm}
\usepackage{setspace}
\setstretch{1.15}
\usepackage{enumitem}
\usepackage{booktabs}
\usepackage{color}
\usepackage{hyperref}
\definecolor{ntr}{rgb}{0.8,0,0}
\hypersetup{colorlinks=true,linkcolor=ntr,citecolor=ntr,urlcolor=black}

% Portada (Luis Cabral): titulo \LARGE sans-serif alineado a la izquierda, sin numero de pagina
\makeatletter
    \renewcommand{\@maketitle}{%
    \newpage\parindent0in\null\vskip2em%
    {\LARGE\textbf{\textsf{\@title}}\par\vskip2em}%
    \@author\par\vskip2em%
    \@date%
    \par
    \thispagestyle{empty}\setcounter{page}{1}%
    }%
\makeatother

\DeclareMathOperator{\Var}{Var}
\DeclareMathOperator{\Cov}{Cov}
\DeclareMathOperator{\Corr}{Corr}
\DeclareMathOperator{\E}{\mathbb{E}}
\DeclareMathOperator{\N}{Normal}
\DeclareMathOperator{\LP}{\mathbb{L}}
\newcommand\independent{\protect\mathpalette{\protect\independenT}{\perp}}
\def\independenT#1#2{\mathrel{\rlap{$#1#2$}\mkern2mu{#1#2}}}
